% sections/02_introduction.tex

\section{Introduction}

Road accidents represent a persistent public safety and urban planning challenge.
Understanding what drives accident severity---and specifically what causes
accidents to have prolonged traffic impact---enables smarter infrastructure
design, better emergency response, and more effective policy.
The US Accidents dataset, compiled by Moosavi et al.\ \cite{moosavi2019dataset},
provides one of the most comprehensive real-world accident records available for
such analysis.

To frame our investigation, we follow the \textbf{CPQH} structure---Context,
Problem, Questions, Hypothesis---a method that grounds our visual storytelling in
facts, real issues, and investigative questions before forming claims.

\begin{table}[H]
\centering
\begin{tabularx}{\linewidth}{@{} l X @{}}
  \toprule
  \textbf{CPQH Stage} & \textbf{Our Application} \\ \midrule
  Context
    & US road accidents 2016--2022 across 49 states; severity scored 1--4
      based on traffic disruption duration. \\[4pt]
  Problem
    & Traffic severity is not random. Environmental and temporal factors may
      systematically increase accident-induced traffic disruption. \\[4pt]
  Questions
    & Does adverse weather prolong accident-induced disruption?
      Do rush hours amplify severity?
      Does road infrastructure moderate severity? \\[4pt]
  Hypotheses
    & H1: Bad weather increases accident severity.
      H2: Rush hours increase accident severity.
      H3 \& H4: Extended hypotheses derived from the severity definition. \\
  \bottomrule
\end{tabularx}
\caption{CPQH framework applied to this project}
\end{table}
